\appendix
\section{Приложение: Вывод уравнения в координатах \texorpdfstring{$(\xi, \eta)$}{(xi, eta)}}	
Уравнение неразрывности в полярных координатах:
\[
\frac{1}{r} \frac{\partial}{\partial r} \left( \rho r \frac{\partial \varphi}{\partial r} \right) + \frac{1}{r^2} \frac{\partial}{\partial \theta} \left( \rho \frac{\partial \varphi}{\partial \theta} \right) = 0.
\]
Переходя к переменным $(\xi, \eta)$ с учетом $\frac{\partial}{\partial r} = e^{-\xi} \frac{\partial}{\partial \xi}$, $\frac{\partial}{\partial \theta} = \frac{\partial}{\partial \eta}$, $r = e^{\xi}$:
\[
\frac{\partial}{\partial \xi} \left( \rho \frac{\partial \varphi}{\partial \xi} \right) + \frac{\partial}{\partial \eta} \left( \rho \frac{\partial \varphi}{\partial \eta} \right) = 0.
\]


\section{Вывод коэффициентов пятиточечного шаблона для ПНМ}

В полностью неявном методе (ПНМ) уравнение неразрывности решается одновременно для всей области. Для этого в каждом узле $(i, j)$ записывается дискретный аналог уравнения в виде невязки $F_{i, j} = 0$.

\subsection{Дискретизация и линеаризация}

Используя центральные разности для потоков (аналогично разделу \ref{sec:slor_derivation}), запишем значение невязки в узле $(i, j)$:
\[
F_{i,j} = \frac{\rho_{i+1/2,j}(\varphi_{i+1,j} - \varphi_{i,j}) - \rho_{i-1/2,j}(\varphi_{i,j} - \varphi_{i-1,j})}{\Delta \xi^2} + \frac{\rho_{i,j+1/2}(\varphi_{i,j+1} - \varphi_{i,j}) - \rho_{i,j-1/2}(\varphi_{i,j} - \varphi_{i,j-1})}{\Delta \eta^2}
\]
Для решения нелинейной системы методом замороженной плотности ищется поправка $\delta\varphi$, такая что:
\[
\frac{\partial F_{i,j}}{\partial \varphi} \delta\varphi = -F_{i,j}
\]
В предположении «замороженной» плотности ($\rho \approx \text{const}$ на шаге итерации), это приводит к пятиточечному уравнению:
\[
\hat{A}_{i,j} \delta \varphi_{i,j-1} + A_{i,j} \delta \varphi_{i-1,j} + B_{i,j} \delta \varphi_{i,j} + C_{i,j} \delta \varphi_{i+1,j} + \hat{C}_{i,j} \delta \varphi_{i,j+1} = -F_{i,j}
\]

\subsection{Коэффициенты матрицы ПНМ}

Сравнивая невязку с общим видом пятиточечной схемы, получаем формулы для коэффициентов:
\begin{itemize}
	\item $A_{i,j} = \displaystyle \frac{\rho_{i-1/2,j}}{\Delta \xi^2}$ — связь с левым узлом ($i-1$).
	\item $C_{i,j} = \displaystyle \frac{\rho_{i+1/2,j}}{\Delta \xi^2}$ — связь с правым узлом ($i+1$).
	\item $\hat{A}_{i,j} = \displaystyle \frac{\rho_{i,j-1/2}}{\Delta \eta^2}$ — связь с нижним узлом ($j-1$).
	\item $\hat{C}_{i,j} = \displaystyle \frac{\rho_{i,j+1/2}}{\Delta \eta^2}$ — связь с верхним узлом ($j+1$).
	\item $B_{i,j} = \displaystyle -\left( \frac{\rho_{i+1/2,j} + \rho_{i-1/2,j}}{\Delta \xi^2} + \frac{\rho_{i,j+1/2} + \rho_{i,j-1/2}}{\Delta \eta^2} \right)$ — центральный коэффициент.
\end{itemize}

Данная система уравнений имеет пятидиагональную структуру и решается в данной работе с помощью сильно неявной процедуры (SIP) Стоуна.

\section{Вывод коэффициентов трехдиагональной системы метода SLOR}
\label{sec:slor_derivation}

Для численного решения уравнения неразрывности:
\[
\frac{\partial}{\partial \xi} \left( \rho \frac{\partial \varphi}{\partial \xi} \right) + \frac{\partial}{\partial \eta} \left( \rho \frac{\partial \varphi}{\partial \eta} \right) = 0
\]
используется аппроксимация центральными разностями второго порядка в узле $(i, j)$.

\subsection{Дискретизация потоков}

Поток через правую грань контрольного объема ($i+1/2, j$):
\[
\left( \rho \frac{\partial \varphi}{\partial \xi} \right)_{i+1/2, j} \approx \rho_{i+1/2, j} \frac{\varphi_{i+1, j} - \varphi_{i, j}}{\Delta \xi}
\]
Поток через левую грань ($i-1/2, j$):
\[
\left( \rho \frac{\partial \varphi}{\partial \xi} \right)_{i-1/2, j} \approx \rho_{i-1/2, j} \frac{\varphi_{i, j} - \varphi_{i-1, j}}{\Delta \xi}
\]
Аналогично для направлений по $\eta$:
\[
\left( \rho \frac{\partial \varphi}{\partial \eta} \right)_{i, j+1/2} \approx \rho_{i, j+1/2} \frac{\varphi_{i, j+1} - \varphi_{i, j}}{\Delta \eta}, \quad \left( \rho \frac{\partial \varphi}{\partial \eta} \right)_{i, j-1/2} \approx \rho_{i, j-1/2} \frac{\varphi_{i, j} - \varphi_{i, j-1}}{\Delta \eta}
\]

\subsection{Формирование системы уравнений}

Подставляя аппроксимации в уравнение и умножая на $-1$, получим:
\[
\frac{\rho_{i-1/2,j}(\varphi_{i,j} - \varphi_{i-1,j}) - \rho_{i+1/2,j}(\varphi_{i+1,j} - \varphi_{i,j})}{\Delta \xi^2} + \frac{\rho_{i,j-1/2}(\varphi_{i,j} - \varphi_{i,j-1}) - \rho_{i,j+1/2}(\varphi_{i,j+1} - \varphi_{i,j})}{\Delta \eta^2} = 0
\]
Группируя слагаемые при неизвестных на текущей строке $j$ ($\varphi_{i-1,j}, \varphi_{i,j}, \varphi_{i+1,j}$):
\[
A_i \varphi_{i-1,j} + B_i \varphi_{i,j} + C_i \varphi_{i+1,j} = D_i
\]
где коэффициенты определяются как:
\begin{itemize}
	\item $A_i = \displaystyle \frac{\rho_{i-1/2,j}}{\Delta \xi^2}$ — влияние предыдущего узла.
	\item $C_i = \displaystyle \frac{\rho_{i+1/2,j}}{\Delta \xi^2}$ — влияние следующего узла.
	\item $B_i = \displaystyle -\left( \frac{\rho_{i+1/2,j} + \rho_{i-1/2,j}}{\Delta \xi^2} + \frac{\rho_{i,j+1/2} + \rho_{i,j-1/2}}{\Delta \eta^2} \right)$ — центральный коэффициент.
	\item $D_i = \displaystyle - \frac{\rho_{i,j+1/2} \varphi_{i,j+1} + \rho_{i,j-1/2} \varphi_{i,j-1}}{\Delta \eta^2}$ — вклад соседних строк.
\end{itemize}

Данный вид системы обеспечивает выполнение принципа максимума и диагональное преобладание ($|B_i| = A_i + C_i + \text{положительные члены}$), что гарантирует устойчивость и сходимость итерационного процесса.
